\scanpage{47}
graphe discret: sans arêtes. Parties discrètes d'un graphe (discrètes pour la structure induite) doublets, triades\inserted{, quadruplets, pentuplets, hextuplets}; parties discrètes de card $2,3$.

\deleted{graphe} graphe \unsure{bi-discret|bidiscret}: graphe \inserted{dont l'ens des sommets} $\Gamma$, décomposé en deux parties discrètes
$A,A'$ $(\Gamma=A\amalg A')$ \deleted{\unsure{avec}} tel que
\deleted{$\forall x\in A\ \exists!x'\in A'$ tel que $\{x,x'\}$}
\inserted{la relation $R$ entre $A,A'$ définie par}
$\{x,x'\}\in R\Longleftrightarrow\{x,x'\}$ est un doublet (i.e. $x,x'$ non liés) est bijective
\unsure{soit|est|sauf} un doublet. Cette décomposition est unique sauf si
$\operatorname{Card} A=2$ \inserted{i.e. $\operatorname{Card}\Gamma=4$} (les graphes de ce dernier type sont \unsure{tous} isom. à
\diagram{\begin{tikzcd}[ampersand replacement=\&,row sep=tiny,column sep=small]
a\arrow[dr,dash]\& b\arrow[dl,dash]\\ a'\& b'
\end{tikzcd}}
$A_2\amalg A_2$, i.e. somme disjointe de deux segments -- et le
choix d'une \deleted{\unsure{A|$\Delta$}} décomposition \inserted{de $\Gamma$} en $A,A'$ équivaut à la donnée d'une bijection entre ces segments.
Bidoublets, bitriades \ldots\ \inserted{biquadruplets, bipentuplets, bihextuples \ldots}
\deleted{parties bi-\unsure{discrètes} d'ordre $n$ (bi}

% slanted
\margin{\uline{NB} les bitriades sont exactement les hexagones}

Dans un graphe cubique, les bases et bibases
sont les \deleted{\unsure{syst\textsuperscript{x}|systèmes} \unsure{parties|$\Gamma=A,A'$}} hextuplets et bihextuplets.

\inserted{Deux parties discrètes \inserted{$A,A'$} d'un graphe sont dites \deleted{``opposées''} ``biliées'' si elles sont disjointes et si leur réunion est \deleted{\unsure{bidiscr.}} un graphe \unsure{(avec ses)} bidiscret pour sa décomposition naturelle.}
% slanted
\margin{elles sont dites liées si elles sont disjointes et si leur réunion est encore discrète. Elles sont dites \uline{opposées} si tt él. de l'une est lié à tt él de l'autre (donc elles sont disjointes)}

\uline{NB} La catégorie des graphes bi-discrets d'ordre $n$ est
l'équivalente canonique à la catégorie $(\mathrm{Ens}_2)\times(\mathrm{Ens})_n$.

Deux triades $\tau=\{a,b,c\}$, $\tau'=\{x,y,z\}$ sont \inserted{donc} \deleted{dites} \uline{liées}
si \deleted{leur réunion} elles sont disjointes et leur réunion
est discrète, (i.e. ssi leur réunion est un hextuplet \unsure{\S} \inserted{ou encore une base.}

\deleted{elles sont dites biliées}

\uline{Proposition} Soit $\Delta=(S,A)$ un graphe cubique, soit
$\mathrm{Tri}_{\Delta}$ l'ens de ses triades, considérons sur $\mathrm{Tri}$, la structure
de graphe définie par la relation des ``triades liées''
précédente. Alors \textcircled{\scriptsize 1} $\mathrm{Tri}$ est \deleted{\unsure{réunion disjointe}} \inserted{un graphe somme} de
\deleted{\unsure{des}} \inserted{\uline{hexagones}} \inserted{(i.e. ses comp. connexes sont des hexagones)}.
\margin{$\operatorname{Card}\mathrm{Tri}_{\Delta}=\unsure{720|{\mathord{?}}20}$}
% slanted
\margin{Hex \unsure{gr.} \deleted{hexagones distingués} des triades: chaque triade appartient à un unique \deleted{hex.} \inserted{hexagone} \deleted{distingué} $\operatorname{Card}\mathrm{Hex}=\unsure{120|720}$}

\textcircled{\scriptsize\unsure{2|1}} Deux triades sont \deleted{biliées} \inserted{opposées au sens plus haut} ssi elles sont sur un $=$ hexagone et \inserted{y} sont \deleted{\unsure{pas liées}} \inserted{opposées} (donc pour tte triade $\exists!$ triade opposée)
\textcircled{\scriptsize 3} Deux triades \inserted{$\tau,\tau'$} sont biliées ssi elles appartiennent à un $=$ \deleted{\unsure{hexagone}} hexagone,
% archive sheet number: 46
